% Background Fields: Position-Dependent Coupling
% Source: docs/background_fields.md (migrated to LaTeX)
% Uses macros from preamble.tex

\section{Background Fields: Position-Dependent Coupling}
\label{sec:background-fields}

Background fields are non-dynamical tensors that appear in the Lagrangian
but are \emph{not} varied in the Euler--Lagrange derivation. They survive as
(possibly position-dependent) coefficients in the equations of motion.

This feature is essential for simulating the Gertsenshtein effect
(Gertsenshtein 1962; Domcke \& Garcia-Cely 2023~\cite{domcke2023simple}), where a static external
magnetic field $\Bzero(\bm{x})$ catalyses photon-graviton conversion. It also enables
general probe-field approximations and theories with fixed external sources.

\subsection{Pipeline Trace}
\label{sec:bg-pipeline-trace}

\subsubsection{1. TOML Declaration}
\label{sec:bg-toml}

\begin{lstlisting}[style=toml]
[[background_fields]]
name = "G"
type = "scalar"
components = ["g0 * Exp[-(x[]^2 + y[]^2) / (2 * R^2)]"]
\end{lstlisting}

Component values can be:

\begin{itemize}
  \item \textbf{Numbers}: \texttt{0}, \texttt{1.0} (constant)
  \item \textbf{Strings}: Wolfram expressions evaluated symbolically (may depend on coordinates)
\end{itemize}

\paragraph{Index Convention: \texttt{components} always specifies covariant (lowered-index) values}
\label{sec:bg-index-convention}

For a vector background field \texttt{A}, \texttt{components[$\mu$]} defines $A_\mu$ (covariant, lower index).

\begin{lstlisting}
components = [A_0, A_1, ..., A_{d-1}]   <- COVARIANT (A with subscript)
\end{lstlisting}

This is the convention throughout TIDAL's Wolfram pipeline:
\begin{itemize}
  \item \texttt{SetBackgroundFieldDownValues} sets \texttt{field[\{$\mu$, -chart\}] = components[$\mu$]} (covariant \texttt{-chart})
  \item \texttt{EvaluatePDBackgroundField} processes covariant \texttt{\{$\mu$, -chart\}} forms
  \item xAct's \texttt{ComponentValue} stores covariant components
\end{itemize}

\textbf{Contravariant forms are derived automatically:}
\begin{itemize}
  \item For \textbf{Minkowski} metrics (\texttt{metric = "minkowski"}): $A^\mu \approx A_\mu$ (spatial components identical; temporal $A^0 = -A_0$, but typically zero)
  \item For \textbf{curved diagonal} metrics (\texttt{metric = "diagonal"}): $A^\mu = g^{\mu\mu} A_\mu = A_\mu / g_{\mu\mu}$, computed via \texttt{Simplify[(A\_$\mu$) / (g\_$\mu\mu$)]}
\end{itemize}

\textbf{Example: Dipolar magnetic field in spherical coordinates} ($\text{diag}[-1, 1, r^2, r^2\sin^2\theta]$):
\begin{lstlisting}
A_theta = Bpeak*z0^3/(2r^2)    <- what you put in components[2]
A^theta = g^{tt} A_theta = A_theta/r^2 = Bpeak*z0^3/(2r^4)   <- computed automatically
\end{lstlisting}

\textbf{Gauge potentials}: Use the covariant gauge potential $A_\mu$ in \texttt{components}.
The Faraday tensor $F_{\mu\nu} = \partial_\mu A_\nu - \partial_\nu A_\mu$ is metric-independent (exterior derivative),
so $F$ is determined correctly from covariant $A_\mu$ alone.

\subsubsection{2. WLS Generation (\texttt{\_derive.py})}
\label{sec:bg-wls-gen}

\begin{itemize}
  \item \texttt{\_wls\_fields()} emits \texttt{DefTensor[G[], M2]} + \texttt{ComponentValue[G[\{0,-Cart\}], expr]}
  \item \textbf{Scalar backgrounds}: Explicit \texttt{ReplaceAll} via \texttt{\_wls\_scalar\_background\_substitution()} before decomposition (\texttt{ToBasis} doesn't trigger \texttt{ComponentValue} for rank-0 tensors)
  \item \textbf{Vector/tensor backgrounds}: Handled by \texttt{ComponentValue} + \texttt{ToBasis} in \texttt{DecomposeToComponents}
\end{itemize}

\subsubsection{3. Wolfram Pipeline}
\label{sec:bg-wolfram}

\begin{itemize}
  \item \texttt{VarD} is called \textbf{only} on dynamical fields --- background fields are correctly \emph{not} varied
  \item The background field's symbolic expression survives as a coefficient in the EOM
\end{itemize}

\subsubsection{4. JSON Export (\texttt{ExportJSON.wl})}
\label{sec:bg-json-export}

\begin{itemize}
  \item \texttt{IsCoordinateDependentCoefficient} detects xCoba symbols in coefficients
  \item Sets \texttt{coordinate\_dependent: ["x", "y"]} and \texttt{coefficient\_symbolic: "g0*Exp[...]"}
  \item \texttt{time\_dependent: true} if the coefficient contains \texttt{t[]}
\end{itemize}

\subsubsection{5. Python PDE Evaluation (\texttt{pde\_builder.py})}
\label{sec:bg-pde-eval}

\begin{itemize}
  \item \texttt{\_mathematica\_to\_python()} converts Mathematica \texttt{InputForm} to Python ($\sim$30 functions)
  \item \texttt{\_resolve\_coefficient\_at\_point()} evaluates on numpy grid arrays via \texttt{eval()} with restricted namespace (\texttt{\_\_builtins\_\_: \{\}})
\end{itemize}

\subsubsection{6. Measurement Module (\texttt{\_energy.py})}
\label{sec:bg-measurement}

\begin{itemize}
  \item \texttt{evaluate\_coefficient()} in \texttt{\_eval\_utils.py} provides standalone evaluation
  \item \texttt{\_resolve\_mass\_squared()} returns \texttt{ndarray} for position-dependent mass terms
  \item \texttt{\_compute\_virial\_potential()} evaluates position-dependent coefficients on the grid
\end{itemize}

\subsection{Caching Architecture}
\label{sec:bg-caching}

Four levels of caching prevent redundant evaluation:

\begin{table}[htbp]
\centering
\caption{Coefficient caching levels in TIDAL's PDE evaluation pipeline.}
\label{tab:bg-caching}
\begin{tabular}{@{}cllp{5cm}@{}}
\toprule
Level & Cache & What & Lifetime \\
\midrule
L0 & \texttt{\_preresolved} & Fully constant coefficients (no \texttt{eval()} needed) & Instance (construction) \\
L1 & \texttt{\_expr\_cache} & Mathematica$\to$Python string conversion & Instance (singleton) \\
L2 & \texttt{\_spatial\_cache} & Grid arrays for spatial-only coefficients & Instance (populated on first RHS call) \\
L3 & \texttt{coeff\_cache} & Per-RHS-call results & Single \texttt{\_compute\_rhs\_for\_component()} call \\
\bottomrule
\end{tabular}
\end{table}

\textbf{Key behavior}: Spatial-only coefficients (\texttt{position\_dependent=True}, \texttt{time\_dependent=False}) are evaluated once on the grid and cached persistently across all timesteps. Time-dependent coefficients are re-evaluated each substep. Fully constant coefficients (neither position- nor time-dependent) are resolved once at construction time (L0) with no \texttt{eval()} at runtime.

\subsection{Supported Patterns}
\label{sec:bg-patterns}

\begin{table}[htbp]
\centering
\caption{Background field patterns supported by TIDAL.}
\label{tab:bg-patterns}
\begin{tabular}{@{}llp{4cm}@{}}
\toprule
Pattern & Example & Status \\
\midrule
Constant background & \texttt{components = ["B0"]} & Fully supported \\
Position-dependent scalar & \texttt{components = ["g0 * Exp[-r\^{}2]"]} & Fully supported \\
Position-dependent vector & \texttt{components = [0, 0, "B0 * Sin[x[]]"]} & Fully supported \\
Time-dependent background & \texttt{components = ["g0 * Cos[t[]]"]} & Works (Python-level tested) \\
Background field gradient & \texttt{CD[-a][G[]]} in Lagrangian & \textbf{NOT supported} --- raises error \\
\bottomrule
\end{tabular}
\end{table}

\subsection{Background Field Validity: Externally Imposed Fields}
\label{sec:bg-validity}

\textbf{TIDAL does NOT validate that background fields satisfy the background equations of
motion.} Validation is purely structural (type, rank, component count, name collisions
--- see \texttt{\_derive\_validate.py}). The user is responsible for ensuring physical consistency.

For linearized perturbation theory, background fields are treated as \textbf{externally imposed}
--- maintained by some external source (e.g., a magnet, stellar interior currents, or an
applied potential). The perturbation equations are valid regardless of whether the
background configuration is self-consistent, because:

\begin{enumerate}
  \item The background is \emph{not} varied in the Euler--Lagrange derivation (\texttt{VarD} acts only on
    dynamical perturbation fields)
  \item The perturbation equations depend on the background only through its values and
    derivatives at each grid point
  \item Whether those values satisfy the background EOM is irrelevant to the linearized dynamics
\end{enumerate}

\textbf{Example}: A spatially-varying magnetic field $\Bzero(z)$ implies $\nabla\times\bm{B} = \bm{J} \neq 0$, violating
vacuum Maxwell equations. This is not a problem --- $\Bzero(z)$ is interpreted as the field
produced by external currents (e.g., a solenoid or stellar magnetosphere). This treatment
is standard in the graviton-photon mixing literature:
\begin{itemize}
  \item Raffelt \& Stodolsky (1988, PRD 37:1237) --- treat $B(z)$ as external
  \item Boccaletti et al.\ (1970, Nuovo Cimento 70B:129) --- arbitrary $B(z)$ profiles
  \item Domcke \& Garcia-Cely (2024, JCAP 05:051)~\cite{domcke2023magnetic} --- inhomogeneous $B(z)$
\end{itemize}

\textbf{When self-consistency matters}: For strong-field or nonlinear regimes (not supported
by TIDAL's linearized framework), the background must satisfy its own EOM. In TIDAL's
linearized regime, this is never required.

\subsection{Energy Measurement with Position-Dependent Coefficients}
\label{sec:bg-energy}

Energy measurements (\texttt{tidal measure --what energy}) correctly handle
spatially-varying coefficients. The Hamiltonian context
(\texttt{\_prepare\_hamiltonian\_context} in \texttt{\_energy.py}) pre-resolves all
position-dependent coefficients as full spatial arrays via
\texttt{evaluate\_coefficient()}, and the per-term evaluation computes
$c(\bm{x}) \cdot f_a(\bm{x}) \cdot f_b(\bm{x}) \cdot \sqrtg$ pointwise before spatial
averaging. This covers:

\begin{itemize}
  \item Position-dependent mass terms (e.g., $m^2(z) \cdot \phi^2$)
  \item Localised coupling regions (e.g., Gaussian $\Bzero(z)$ in Gertsenshtein)
  \item Potential wells
\end{itemize}

\subsection{Periodic BC Requirement for Position-Dependent Coefficients}
\label{sec:bg-periodic-bc}

When using periodic BCs with position-dependent coefficients, the
coefficients must be approximately continuous at the domain boundaries.
Non-periodic coefficients (e.g., $B(x) = B_\mathrm{peak}/x^3$ on $[5, 80]$)
break the integration-by-parts identity and cause $O(1)$ energy
non-conservation \textbf{independent of grid resolution and solver tolerance}.

A runtime check (\texttt{check\_periodic\_coefficient\_continuity}) runs
automatically at solver initialisation. It uses a \emph{leak metric} =
$(\text{rel\_jump}) \times (\text{boundary\_fraction})$ to estimate the actual energy leak.
See \texttt{docs/troubleshooting.md} for solutions when this warning fires.

\subsection{Known Limitations}
\label{sec:bg-limitations}

\begin{enumerate}
  \item \textbf{No background field gradients}: \texttt{CD[-a][G[]]} in the Lagrangian will raise
    \texttt{ValueError} at TOML validation time. The scalar \texttt{ReplaceAll} substitution
    doesn't handle derivatives of background fields.

  \item \textbf{Time-dependent backgrounds less tested}: The Python evaluation pipeline
    handles time-dependent expressions correctly, but the Wolfram-side behavior
    (xAct handling time derivatives of time-dependent background expressions)
    is not end-to-end tested with wolframscript.

  \item \textbf{Sign change diagnostic}: If a position-dependent mass term changes sign
    across the grid, a \texttt{UserWarning} is emitted (possible tachyonic instability).
    This is informational, not a hard error.
\end{enumerate}

\subsection{Adding New Mathematica Functions}
\label{sec:bg-new-functions}

To support a new Mathematica function (e.g., \texttt{Gamma[x]}), update
\texttt{\_eval\_utils.py} only --- \texttt{pde\_builder.py} delegates to it automatically:

\begin{enumerate}
  \item \textbf{String conversion}: Add to \texttt{\_FUNCTION\_MAP} in \texttt{\_eval\_utils.py}:

\begin{lstlisting}[style=python]
("Gamma", "gamma"),
\end{lstlisting}

  \item \textbf{Evaluation namespace}: Add to \texttt{build\_eval\_namespace()} in \texttt{\_eval\_utils.py}:

\begin{lstlisting}[style=python]
ns["gamma"] = special.gamma
\end{lstlisting}

  \item \textbf{Test}: Add a conversion test in \texttt{test\_background\_fields.py}.
\end{enumerate}

\texttt{PDEFromSpec.\_mathematica\_to\_python()} and \texttt{PDEFromSpec.\_build\_base\_namespace()}
both delegate to \texttt{\_eval\_utils.py}, so no changes are needed in \texttt{pde\_builder.py}.

\subsection{References}
\label{sec:bg-references}

\begin{itemize}
  \item Gertsenshtein (1962), ``Wave resonance of light and gravitational waves'', JETP 14, 84~\cite{gertsenshtein1962wave} --- original prediction requiring external $B$-field as background
  \item Domcke \& Garcia-Cely (2023), ``A simple derivation of the Gertsenshtein effect'', arXiv:2301.02072~\cite{domcke2023simple} --- modern derivation with inhomogeneous background profiles
  \item Mart\'in-Garc\'ia et al., ``xAct: Efficient tensor computer algebra for Mathematica'' --- symbolic tensor algebra (\texttt{VarD}, \texttt{ComponentValue}, \texttt{ToBasis})
  \item Zwicker (2020), ``py-pde: A Python package for solving partial differential equations'', JOSS 5(48), 2158 --- original PDE backend; FD stencil conventions retained in TIDAL's native operators
\end{itemize}

See \texttt{docs/references.md} for the full citation list.
